\chapter{原子・分子・結晶}
\section{背景：状態・構造・物性を三段階でつなぐ}
原子の電子状態を記述し、複数原子がどのように結合するかを考え、並んだ結晶の物性へ進む。電子配置から導電性までを一つの暗記表にすると、共有結合なら導電性が高い、胞の原子数なら配位数と同じ、といった取り違えが生じる。物理の量子状態と幾何を使い、それぞれの定義へ戻って判断する。

\section{空間軌道と電子状態}\label{theory:atomic:orbital}
非相対論的な中心力近似で、空間軌道は主量子数\(n\)、方位量子数\(l\)、磁気量子数\(m_l\)で表す。
\[
 n=1,2,\ldots,\qquad l=0,\ldots,n-1,\qquad m_l=-l,\ldots,l.
\]
\(l=0,1,2,3\)をs,p,d,fと呼ぶ。電子状態にはさらに\(m_s=\pm1/2\)が必要であり、水素原子の電子にもスピンはある。空間軌道と、そこに置く電子の状態は同じ対象ではない\cite{atomic_quantum}。

Pauli原理により同じ空間軌道の電子は最大2個で逆スピンである。従ってs副殻は2、pは6、dは10、fは14電子まで入る。中性酸素は8電子で\(1s^2 2s^2 2p^4\)。総電子数が8でも\(2s^3\)や\(2s^4\)は容量を超える。原子番号と電子数が等しいのは中性原子で、イオンでは電子増減を数える。

通常の基底状態の構築順では3pの後に4sが入るが、全原子・励起状態・イオンで一定の軌道エネルギー順があるとはしない。Cuは\([\mathrm{Ar}]3d^{10}4s^1\)、Znは\([\mathrm{Ar}]3d^{10}4s^2\)である\cite{atomic_nist2019}。

\begin{tcolorbox}[title=公式正答と独立判断の相違を残す]
令和元年度問14の公式は1、独立初回は4。原典の4は水素の空間軌道と多電子原子のスピンを対比し、電子状態の意味で読むと水素のスピンも必要になる。一方、1に含まれる\(\phi\)を「存在しない」と断言すると、2014年のthorocene研究で報告された\(\phi\)軌道相互作用と整合しない\cite{atomic_phi2014}。多中心の軌道相互作用と純粋な二原子\(\phi\)結合を同一視もしない。公式の採点を保持し、この不一致を未解決の査読対象とする。
\end{tcolorbox}

\section{結合次数と電子対配置}\label{theory:atomic:bonding}
分子軌道の結合次数は\((N_b-N_a)/2\)で、結合性だけでなく反結合性の電子も数える。代表的な\(\mathrm{N_2,O_2,F_2}\)は3,2,1であり、順に共有電子が増え結合が強くなるという説明は逆である。個別分子の結合長・エネルギーは分子軌道占有と合わせて読む。

VSEPRは中心原子まわりの電子領域を数える近似である。水は結合対2、孤立対2のAB2E2。4電子領域の配置と分子形を区別し、CH4は四面体、NH3は三角錐、H2Oは折れ線となる。sp3という共通モデルは、三分子の全原子が同じ四面体形という意味ではない。d関与の形式的混成表記も、全超原子価分子の実在結合をd軌道占有だけで説明する保証ではない。

\section{分子間力：ポテンシャルを微分して力を得る}\label{theory:atomic:interactions}
無極性分子でも瞬間的な電子密度の偏りが隣の分子に双極子を誘起し、分散力が生じる。典型的な水素結合では、電気陰性度の高い原子に結合したHと別の電子対を持つ原子が供与体・受容体として相互作用する。同種分子に限らない\cite{atomic_nonbond}。対称なCH4は永久双極子を持たず、強い分子極性を前提とする水素結合の説明にはならない。

12--6 Lennard--Jonesモデルは
\[
 U(r)=4\varepsilon\left[(\sigma/r)^{12}-(\sigma/r)^6\right],\qquad
 F_r=-\frac{dU}{dr}=24\varepsilon\left[\frac{2\sigma^{12}}{r^{13}}-\frac{\sigma^6}{r^7}\right].
\]
ポテンシャルの反発は正の12乗項、引力は負の6乗項で、力の指数は13と7へ変わる\cite{atomic_lj}。平成29年度問11では反発・引力の逆転を説明できるが、Wの凝集エネルギーを配位数だけへ帰する別候補にも疑義が残る。金属の多電子状態を一つの原子対のモデルへ無条件還元しない。

\section{格子・基底・最近接距離}\label{theory:atomic:geometry}
格子は等価な点を並進で繰り返す数学的集合、基底は各格子点へ付く原子群である。3次元Bravais格子は14種類である\cite{atomic_bravais_ref}。慣用立方胞のFCCは原子数4、最近接配位数12、BCCは原子数2、最近接配位数8。同じ数ではない\cite{openstax_crystals}。

FCCの隅から面中心へのベクトルは\((a/2,a/2,0)\)、BCCの隅から体中心へは\((a/2,a/2,a/2)\)だから
\[
 d_{\mathrm{FCC}}=a/\sqrt2,\qquad d_{\mathrm{BCC}}=\sqrt3a/2.
\]
過去問のdが最近接距離なのか格子定数なのかを各回で確認する。FCCでdからaなら\(a=\sqrt2d\)。BCCでaからdなら\(d=\sqrt3a/2\)。原子半径をrとし接触硬球を置く場合はd=2rである。

HCPの慣用胞規約・理想軸比・すべり条件は第\ref{theory:material:crystals}節と合わせて読む。元素例の構造は温度・圧力・組成に依存する。純Feの通常圧の912℃構造変態と、770℃付近の磁気転移を同一視しない。アモルファスは長距離並進秩序を持たないが短距離の局所秩序はあり得る。

逆格子の原点まわりでWigner--Seitz胞を作った領域が第一Brillouin域である。実格子の胞と逆格子の胞を区別する。diamondはFCCに二原子基底を付け、第一近接4、第二近接12を持つ。

\section{単位胞の密度：分子量と長さの単位をそろえる}\label{theory:atomic:density}
慣用立方胞の原子数z、モル質量M、格子定数aなら
\[
 \rho=\frac{zM}{N_Aa^3}.
\]
胞質量と胞体積の比であり、配位数をzへ代入しない。平成26年度はBCCのz=2、M=52g/mol、a=0.29nm=\(2.9\times10^{-8}\)cmから7.11g/cm3。令和4年度はFCCのz=4、M=0.197kg/mol、a=\(4.1\times10^{-10}\)mから\(1.9056\times10^4\)kg/m3。指定\(N_A=6.0\times10^{23}\)を使用する。OCRの指数33は原典画像の23へ校正し、未校正原文も保存する。

\section{NaClとCsCl：異種の中心はBCC格子点ではない}\label{theory:atomic:ionic}
NaCl型は二つのFCC副格子がずれた構造で異種の最近接配位数6。CsCl型は単純立方格子に異種イオンの基底が付いて配位数8となる。胞の中心と隅が違う元素なので、全点が同等なBCC Bravais格子ではない\cite{openstax_crystals}。Madelung定数はイオンの並びを反映する格子和で、同じ1:1組成というだけで同値とはならない。

NaClに接触硬球のイオン半径\(r_+,r_-\)を置くなら\(a/2=r_++r_-\)。同種の最近接距離\(a/\sqrt2\)は\(\sqrt2(r_++r_-)\)。前半は格子幾何、後半の半径和はモデル仮定である。電子雲に固い表面があり全隣接が厳密に触れるという意味ではない。

\section{炭素の構造と電子状態}\label{theory:atomic:carbon}
diamondは四面体sp3の強い結合網を持ち、通常純結晶は大きなバンドギャップのため良導体ではない\cite{atomic_diamond}。高硬度と高導電性は同じ物性ではない。黒鉛は層内の六角網面と比較的弱い層間結合を持つ。CNTは網面を巻いた単層又は多層の管で、径と巻き方が電子構造へ影響する\cite{atomic_cnt}。

ホウ素ドープしたdiamondやCNTの導電性の例がある\cite{atomic_cnt}。通常純結晶の比較と、ドープ・欠陥・接触条件で変えた材料を分ける。Mohsの最高位10も全人工材料・全方向・全荷重の絶対的順位ではない。

\section{格子振動と磁区応答}\label{theory:atomic:phonon}
格子の正規振動を量子化した励起がphononで、Einsteinモデルは共通の振動数の独立振動子、Debyeモデルは振動数分布を考える。高温の格子比熱はモル当たり3Rへ近づき、低温では励起が減り、3次元Debye近似で格子比熱は\(T^3\)に比例する\cite{atomic_phonon}。金属全体の比熱には電子の寄与もある。

Debye--Waller因子は原子位置の熱変位等による散乱強度の減衰を表す。一方、Barkhausen効果は磁区壁のピン止めと不連続な磁化応答が中心である\cite{atomic_barkhausen}。格子欠陥が磁区壁へ作用し得ることを認めながら、過去問の相対比較では磁区応答を選ぶ。

\section{分子光吸収・光反応と線幅}\label{theory:atomic:optical}
IR振動吸収には双極子の変化、Raman散乱には分極率の変化が関係する\cite{atomic_ir_raman}。CO2の温室効果の主な振動吸収は熱赤外域であり、可視光の振動吸収という説明は適切でない\cite{atomic_ir_co2}。光子のエネルギーとどの自由度を励起するかを先に確認する。

視物質の11cisレチナールの光異性化は視覚の起点になる\cite{atomic_retinal}。光合成ではクロロフィルが光を吸収し電子励起を生じ、NADP+は電子を受け取りNADPHへ還元される役割を持つ\cite{atomic_photosynthesis}。二つを同じ一次光捕集体としない。光素子と超伝導の比較は第\ref{theory:modern:devices}節を参照する。

許容遷移の選択則は線幅が不変という意味ではない。速度分布によるDoppler幅と衝突による幅が存在し、NISTの分子吸収解析でも温度・分子質量・圧力による幅を扱う\cite{atomic_linewidth}。

\section{溶解度と同位体：定義の分母と基準}\label{theory:atomic:solubility}\label{theory:atomic:isotopes}
平成27年度問11の溶解度は100gの飽和溶液中の溶質量である。60℃の水は47.8g、20℃の溶質/水は24/76だから、溶存量15.0947g、析出量37.1053gとなる。100g水当たりの定義へ読み替えず、水の蒸発と結晶への水取込みを無視する条件を明示する。

同位体は同じZでNが違う核種である。単体の構造の違いは同素体、同じAは同重体、周期表の族は別分類。自然試料のδは第\ref{theory:isotope:delta}節の基準比へ戻る。海水硫酸由来NBS-127の\(\delta^{34}\mathrm S\)は約+20.3‰であり、旧CDT脚注もある\cite{atomic_sulfate}。全自然海水を厳密に同値とせず、基準尺度と代表値を区別する。

\section{学習の順序と典型的な誤り}
まず空間軌道の容量と電子状態を区別し、次に電子対配置と結合次数、分子間力へ進む。結晶では格子点と基底を分けて距離を幾何から導き、その後胞質量と体積から密度を求める。導電性・硬度・比熱は電子又は格子のどの自由度で説明するかを確認する。疑義のある全称・不一致を公式正答の番号だけで消さない。

\section{過去問と全候補の検討}
\section{結合次数と電子対配置のモデル}

\subsection{令和元年度 問17}
\label{q:r01-reg-applied-17}

\noindent\textbf{原典：} sources/R01/17\_applied\_science.pdf、PDF物理5頁。

\textbf{公式正答：} 1（single）。

\par\textbf{検討未確定：} Arを分子性結晶と呼ぶ広義分類と単原子性を区別。水の常温液体は圧力等を要する。

\paragraph{問題の要旨} 結合次数・混成モデル・分子間相互作用を比較する。

\Needspace{4\baselineskip}\paragraph{確認した条件・記述（独自要約）}\mbox{}\par

\noindent\textbf{判断対象：} 結合次数・混成モデル・分子間相互作用を比較する。\par

\noindent\textbf{原典確認：} 全5候補・指数・否定・記号を保存原典22頁画像と照合。\par

\paragraph{問われる核心} 結合次数・混成モデル・分子間相互作用を比較する。

\paragraph{必要知識} N2,O2,F2の結合次数は3,2,1で、共有する結合電子対が増えるという1は逆。

\paragraph{考え方・正答を選ぶ理由（査読待ち初稿）} N2,O2,F2の結合次数は3,2,1で、共有する結合電子対が増えるという1は逆。sp3の電子対配置と実際の分子形は違い、NH3やH2Oを正四面体分子と呼ばない。Arの単原子固体を分子性の分類へ含める用語と、通常常圧で水が液体という状態条件を分ける。

\Needspace{4\baselineskip}\paragraph{選択肢の検討（査読待ち初稿）}\mbox{}\par

\noindent\textbf{1：選ぶ候補} N2/O2/F2順に共有電子増で強い。結合次数は3/2/1で逆。反結合性電子の占有も数える。\par

\noindent\textbf{2：選ばない候補} CH4/NH3/H2Oはsp3を反映。4電子領域のモデル。分子形は四面体/三角錐/折れ線と異なる。\par

\noindent\textbf{3：選ばない候補} NaClは静電引力と立方胞。NaCl型の慣用立方胞。\par

\noindent\textbf{4：選ばない候補} 水はHの正電荷と水素結合で常温液体。通常の常圧条件で分子間水素結合が液体を安定化。\par

\noindent\textbf{5：選ばない候補} Ar/ナフタレンは分子性結晶。Arは単原子の分散力固体を広い分子性分類で含める用語。Ar分子や常温固体を意味しない。\par

\paragraph{再利用する考え方} 軌道と電子状態、格子と基底、単位胞原子数と配位数、定義の分母を区別して候補を全て評価する。

\paragraph{共通理論} \hyperref[theory:atomic:bonding]{第\ref*{theory:atomic:bonding}節}

\par 参考：\cite{atomic_nonbond}。

\par 難易度 2（校準初稿）、優先度未評価。

\section{Bravais格子・基底と逆格子}

\subsection{令和7年度 問8}
\label{q:r07-reg-applied-08}

\noindent\textbf{原典：} sources/R07/17\_applied\_science.pdf、PDF物理3頁。

\textbf{公式正答：} 2（single）。

\paragraph{問題の要旨} 格子の並進対称性・Bravais種類と基底を区別する。

\Needspace{4\baselineskip}\paragraph{確認した条件・記述（独自要約）}\mbox{}\par

\noindent\textbf{判断対象：} 格子の並進対称性・Bravais種類と基底を区別する。\par

\noindent\textbf{原典確認：} 全5候補・指数・否定・記号を保存原典22頁画像と照合。\par

\paragraph{問われる核心} 格子の並進対称性・Bravais種類と基底を区別する。

\paragraph{必要知識} 3次元Bravais格子は14種類なので8とする2が不適切。

\paragraph{考え方・正答を選ぶ理由（査読待ち初稿）} 3次元Bravais格子は14種類なので8とする2が不適切。結晶は格子と基底で構成され、NaClの異種イオンやdiamondの二原子基底を格子点の種類として別カウントしない。逆格子で原点まわりのWigner–Seitz胞を第一Brillouin域と呼ぶ。胞原子数と第一/第二配位数は別。

\Needspace{4\baselineskip}\paragraph{選択肢の検討（査読待ち初稿）}\mbox{}\par

\noindent\textbf{1：選ばない候補} 3次元格子は並進対称。格子点は整数並進で互いに等価。\par

\noindent\textbf{2：選ぶ候補} 3次元Bravaisは8種。3次元Bravais格子は14種類。\par

\noindent\textbf{3：選ばない候補} 逆格子Wigner–Seitz胞はBrillouin域。逆格子点まわりのWigner–Seitz胞は第一Brillouin域。\par

\noindent\textbf{4：選ばない候補} diamondの隣接は4と12。四面体第一近接4と同じ副格子の第二近接12。\par

\noindent\textbf{5：選ばない候補} NaClはFCCに属す。FCC Bravais格子にNa/Clの基底が付く。\par

\paragraph{再利用する考え方} 軌道と電子状態、格子と基底、単位胞原子数と配位数、定義の分母を区別して候補を全て評価する。

\paragraph{共通理論} \hyperref[theory:atomic:geometry]{第\ref*{theory:atomic:geometry}節}

\par 参考：\cite{atomic_bravais_ref,openstax_crystals}。

\par 難易度 2（校準初稿）、優先度未評価。

\section{単位胞質量・体積と密度}

\subsection{平成26年度 問14}
\label{q:h26-reg-applied-14}

\noindent\textbf{原典：} sources/H26/17\_applied\_science.pdf、PDF物理6頁。

\textbf{公式正答：} 2（single）。

\paragraph{問題の要旨} BCCの胞原子数と体積から密度を求める。

\Needspace{4\baselineskip}\paragraph{確認した条件・記述（独自要約）}\mbox{}\par

\noindent\textbf{構造：} BCC\par

\noindent\textbf{格子定数：} 0.29nm\par

\noindent\textbf{原子量：} 52\par

\noindent\textbf{指定NA：} 6.0×10\textasciicircum{}23mol\textasciicircum{}-1\par

\noindent\textbf{求める単位：} g/cm3\par

\paragraph{問われる核心} BCCの胞原子数と体積から密度を求める。

\paragraph{必要知識} BCC慣用立方胞では隅8×1/8と体中心1を合わせ2原子。

\paragraph{考え方・正答を選ぶ理由（査読待ち初稿）} BCC慣用立方胞では隅8×1/8と体中心1を合わせ2原子。0.29nm=\allowbreak{}2.9×10\textasciicircum{}-8cmなので、ρ=(2×52g/mol)/(6.0×10\textasciicircum{}23mol\textasciicircum{}-1×(2.9×10\textasciicircum{}-8cm)\textasciicircum{}3)=\allowbreak{}7.11g/cm3、2。長さを立方する前に単位を変える。現在の定義定数へ置き換えず原典指定値を使う。

\Needspace{4\baselineskip}\paragraph{選択肢の検討（査読待ち初稿）}\mbox{}\par

\noindent\textbf{1：選ばない候補} 3.6g/cm3。胞原子数を1とした約半分の値。\par

\noindent\textbf{2：選ぶ候補} 7.1g/cm3。2×52/(6.0×10\textasciicircum{}23×(2.9×10\textasciicircum{}-8)\textasciicircum{}3)=約7.11。\par

\noindent\textbf{3：選ばない候補} 7.9g/cm3。指定量からの計算値より大きい。\par

\noindent\textbf{4：選ばない候補} 10.1g/cm3。指定量からの計算値より大きい。\par

\noindent\textbf{5：選ばない候補} 12.3g/cm3。指定量からの計算値より大きい。\par

\paragraph{再利用する考え方} 軌道と電子状態、格子と基底、単位胞原子数と配位数、定義の分母を区別して候補を全て評価する。

\paragraph{共通理論} \hyperref[theory:atomic:density]{第\ref*{theory:atomic:density}節}

\par 参考：\cite{openstax_crystals}。

\paragraph{関連問題} \hyperref[q:r04-reg-applied-12]{R04-REG-APPLIED-12}

\par 難易度 3（校準初稿）、優先度未評価。

\subsection{令和4年度 問12}
\label{q:r04-reg-applied-12}

\noindent\textbf{原典：} sources/R04/17\_applied\_science.pdf、PDF物理4頁。

\textbf{公式正答：} 4（single）。

\paragraph{問題の要旨} FCCの胞質量をSI体積で割り密度を求める。

\Needspace{4\baselineskip}\paragraph{確認した条件・記述（独自要約）}\mbox{}\par

\noindent\textbf{構造：} FCC\par

\noindent\textbf{格子定数：} 0.41nm\par

\noindent\textbf{原子量：} 197\par

\noindent\textbf{指定NA：} 6.0×10\textasciicircum{}23mol\textasciicircum{}-1\par

\noindent\textbf{求める単位：} kg/m3\par

\paragraph{問われる核心} FCCの胞質量をSI体積で割り密度を求める。

\paragraph{必要知識} FCC慣用胞の原子数は8×1/8+6×1/2=\allowbreak{}4。

\paragraph{考え方・正答を選ぶ理由（査読待ち初稿）} FCC慣用胞の原子数は8×1/8+6×1/2=\allowbreak{}4。胞質量は4×0.197/(6.0×10\textasciicircum{}23)kg、体積は(4.1×10\textasciicircum{}-10)\textasciicircum{}3m3。ρ=1.9056×10\textasciicircum{}4kg/m3なので4。197g/molをkg/molへ直す。H26と同じ方法だが構造・単位・数値が異なる。

\Needspace{4\baselineskip}\paragraph{選択肢の検討（査読待ち初稿）}\mbox{}\par

\noindent\textbf{1：選ばない候補} 7.9×10\textasciicircum{}3kg/m3。計算密度に達しない。\par

\noindent\textbf{2：選ばない候補} 1.2×10\textasciicircum{}4kg/m3。計算密度より小さい。\par

\noindent\textbf{3：選ばない候補} 1.5×10\textasciicircum{}4kg/m3。計算密度より小さい。\par

\noindent\textbf{4：選ぶ候補} 1.9×10\textasciicircum{}4kg/m3。胞原子数4とa=\allowbreak{}4.1×10\textasciicircum{}-10mから約1.91×10\textasciicircum{}4。\par

\noindent\textbf{5：選ばない候補} 2.3×10\textasciicircum{}4kg/m3。計算密度より大きい。\par

\paragraph{再利用する考え方} 軌道と電子状態、格子と基底、単位胞原子数と配位数、定義の分母を区別して候補を全て評価する。

\paragraph{共通理論} \hyperref[theory:atomic:density]{第\ref*{theory:atomic:density}節}

\par 参考：\cite{openstax_crystals}。

\paragraph{関連問題} \hyperref[q:h26-reg-applied-14]{H26-REG-APPLIED-14}

\par 難易度 3（校準初稿）、優先度未評価。

\section{格子・配位数と最近接距離}

\subsection{平成24年度 問5}
\label{q:h24-reg-applied-05}

\noindent\textbf{原典：} sources/H24/17\_applied\_science.pdf、PDF物理2頁。

\textbf{公式正答：} 2（single）。

\paragraph{問題の要旨} FCCの最近接距離から格子定数を求める。

\Needspace{4\baselineskip}\paragraph{確認した条件・記述（独自要約）}\mbox{}\par

\noindent\textbf{判断対象：} FCCの最近接距離から格子定数を求める。\par

\noindent\textbf{原典確認：} 全5候補・指数・否定・記号を保存原典22頁画像と照合。\par

\paragraph{問われる核心} FCCの最近接距離から格子定数を求める。

\paragraph{必要知識} FCCでは面対角線上の隅と面中心が最近接で、d=√2a/2。

\paragraph{考え方・正答を選ぶ理由（査読待ち初稿）} FCCでは面対角線上の隅と面中心が最近接で、d=√2a/2。従ってa=√2dで2。面中心と隅のベクトル(a/2,a/2,0)の長さからも導ける。単位胞の原子数4と配位数12を、この距離式の係数として代入しない。

\Needspace{4\baselineskip}\paragraph{選択肢の検討（査読待ち初稿）}\mbox{}\par

\noindent\textbf{1：選ばない候補} d。FCCの最近接距離は胞辺より短い。\par

\noindent\textbf{2：選ぶ候補} √2d。面対角線を半分にした距離d=\allowbreak{}a/√2からa=√2d。\par

\noindent\textbf{3：選ばない候補} d/√2。距離と格子定数を逆にした値。\par

\noindent\textbf{4：選ばない候補} 2d/√3。BCCなら現れる関係でFCCではない。\par

\noindent\textbf{5：選ばない候補} √3d。FCCの面対角線の幾何に合わない。\par

\paragraph{再利用する考え方} 軌道と電子状態、格子と基底、単位胞原子数と配位数、定義の分母を区別して候補を全て評価する。

\paragraph{共通理論} \hyperref[theory:atomic:geometry]{第\ref*{theory:atomic:geometry}節}

\par 参考：\cite{openstax_crystals}。

\paragraph{関連問題} \hyperref[q:h28-reg-applied-07]{H28-REG-APPLIED-07}, \hyperref[q:r01-ret-applied-07]{R01-RET-APPLIED-07}

\par 難易度 3（校準初稿）、優先度未評価。

\subsection{平成28年度 問7}
\label{q:h28-reg-applied-07}

\noindent\textbf{原典：} sources/H28/17\_applied\_science.pdf、PDF物理3頁。

\textbf{公式正答：} 2（single）。

\paragraph{問題の要旨} FCCの最近接距離から格子定数を求める。

\Needspace{4\baselineskip}\paragraph{確認した条件・記述（独自要約）}\mbox{}\par

\noindent\textbf{判断対象：} FCCの最近接距離から格子定数を求める。\par

\noindent\textbf{原典確認：} 全5候補・指数・否定・記号を保存原典22頁画像と照合。\par

\paragraph{問われる核心} FCCの最近接距離から格子定数を求める。

\paragraph{必要知識} FCCでは面対角線上の隅と面中心が最近接で、d=√2a/2。

\paragraph{考え方・正答を選ぶ理由（査読待ち初稿）} FCCでは面対角線上の隅と面中心が最近接で、d=√2a/2。従ってa=√2dで2。面中心と隅のベクトル(a/2,a/2,0)の長さからも導ける。単位胞の原子数4と配位数12を、この距離式の係数として代入しない。

\Needspace{4\baselineskip}\paragraph{選択肢の検討（査読待ち初稿）}\mbox{}\par

\noindent\textbf{1：選ばない候補} d。FCCの最近接距離は胞辺より短い。\par

\noindent\textbf{2：選ぶ候補} √2d。面対角線を半分にした距離d=\allowbreak{}a/√2からa=√2d。\par

\noindent\textbf{3：選ばない候補} d/√2。距離と格子定数を逆にした値。\par

\noindent\textbf{4：選ばない候補} 2d/√3。BCCなら現れる関係でFCCではない。\par

\noindent\textbf{5：選ばない候補} √3d。FCCの面対角線の幾何に合わない。\par

\paragraph{再利用する考え方} 軌道と電子状態、格子と基底、単位胞原子数と配位数、定義の分母を区別して候補を全て評価する。

\paragraph{共通理論} \hyperref[theory:atomic:geometry]{第\ref*{theory:atomic:geometry}節}

\par 参考：\cite{openstax_crystals}。

\paragraph{関連問題} \hyperref[q:h24-reg-applied-05]{H24-REG-APPLIED-05}

\par 難易度 3（校準初稿）、優先度未評価。

\subsection{令和元年度再試験 問7}
\label{q:r01-ret-applied-07}

\noindent\textbf{原典：} sources/R01\_retest/17\_applied\_science.pdf、PDF物理3頁。

\textbf{公式正答：} 1（single）。

\paragraph{問題の要旨} BCCの格子定数から最近接距離を求める。

\Needspace{4\baselineskip}\paragraph{確認した条件・記述（独自要約）}\mbox{}\par

\noindent\textbf{判断対象：} BCCの格子定数から最近接距離を求める。\par

\noindent\textbf{原典確認：} 全5候補・指数・否定・記号を保存原典22頁画像と照合。\par

\paragraph{問われる核心} BCCの格子定数から最近接距離を求める。

\paragraph{必要知識} BCCの体中心(d/2,d/2,d/2)と隅(0,0,0)の距離は√3d/2で1。

\paragraph{考え方・正答を選ぶ理由（査読待ち初稿）} BCCの体中心(d/2,d/2,d/2)と隅(0,0,0)の距離は√3d/2で1。FCCの面対角線と混同しない。ここで原典のdは格子定数であり、H24等のd=最近接距離と記号の役割が違う。

\Needspace{4\baselineskip}\paragraph{選択肢の検討（査読待ち初稿）}\mbox{}\par

\noindent\textbf{1：選ぶ候補} √3d/2。体対角線の隅から体中心までの半分。\par

\noindent\textbf{2：選ばない候補} d/√2。FCCの関係。\par

\noindent\textbf{3：選ばない候補} √3d。体対角線全長で最近接の2倍。\par

\noindent\textbf{4：選ばない候補} √2d。面対角線全長。\par

\noindent\textbf{5：選ばない候補} d。隅どうしの胞辺で体中心はさらに近い。\par

\paragraph{再利用する考え方} 軌道と電子状態、格子と基底、単位胞原子数と配位数、定義の分母を区別して候補を全て評価する。

\paragraph{共通理論} \hyperref[theory:atomic:geometry]{第\ref*{theory:atomic:geometry}節}

\par 参考：\cite{openstax_crystals}。

\paragraph{関連問題} \hyperref[q:h24-reg-applied-05]{H24-REG-APPLIED-05}

\par 難易度 3（校準初稿）、優先度未評価。

\subsection{令和3年度 問8}
\label{q:r03-reg-applied-08}

\noindent\textbf{原典：} sources/R03/17\_applied\_science.pdf、PDF物理3頁。

\textbf{公式正答：} 2（single）。

\par\textbf{検討未確定：} 元素例は通常相を想定。HCPを全元素で理想最密軸比とする読みは避け、出題時全相条件は未指定。

\paragraph{問題の要旨} 結晶の配位数と塑性すべり条件を区別する。

\Needspace{4\baselineskip}\paragraph{確認した条件・記述（独自要約）}\mbox{}\par

\noindent\textbf{判断対象：} 結晶の配位数と塑性すべり条件を区別する。\par

\noindent\textbf{原典確認：} 全5候補・指数・否定・記号を保存原典22頁画像と照合。\par

\paragraph{問われる核心} 結晶の配位数と塑性すべり条件を区別する。

\paragraph{必要知識} FCCの配位数12を8とした2が不適切。

\paragraph{考え方・正答を選ぶ理由（査読待ち初稿）} FCCの配位数12を8とした2が不適切。BCCは8で充填68\%。胞原子数BCC2/FCC4と最近接配位数8/12を混同しない。すべりは結晶面と方向へ分解した剪断応力で判断する。HCP元素例を全温度圧力や理想軸比の保証へ広げない。

\Needspace{4\baselineskip}\paragraph{選択肢の検討（査読待ち初稿）}\mbox{}\par

\noindent\textbf{1：選ばない候補} αFe/Cr/MoはBCCで充填約68\%。BCC硬球充填率π√3/8≈0.680。\par

\noindent\textbf{2：選ぶ候補} Al/Ni/CuはFCCで配位数8。FCC最近接は12で8はBCC。\par

\noindent\textbf{3：選ばない候補} Mg/Zn/Cd/αTiはHCP。通常相の六方配列。Zn等の軸比を理想最密値へ固定しない。\par

\noindent\textbf{4：選ばない候補} 分解剪断応力が臨界に達しすべる。τ=σcosφcosλのSchmid則に対応。\par

\noindent\textbf{5：選ばない候補} 温度等で構造変化は同素変態。同元素の結晶構造が相条件で変わる。\par

\paragraph{再利用する考え方} 軌道と電子状態、格子と基底、単位胞原子数と配位数、定義の分母を区別して候補を全て評価する。

\paragraph{共通理論} \hyperref[theory:atomic:geometry]{第\ref*{theory:atomic:geometry}節}, \hyperref[theory:material:crystals]{第\ref*{theory:material:crystals}節}

\par 参考：\cite{openstax_crystals}。

\paragraph{関連問題} \hyperref[q:r04-reg-applied-08]{R04-REG-APPLIED-08}

\par 難易度 2（校準初稿）、優先度未評価。

\subsection{令和4年度 問8}
\label{q:r04-reg-applied-08}

\noindent\textbf{原典：} sources/R04/17\_applied\_science.pdf、PDF物理3頁。

\textbf{公式正答：} 5（single）。

\par\textbf{検討未確定：} 元素例の全相境界・出題時測定条件は未指定。912℃は通常圧純Feの代表構造変態値。

\paragraph{問題の要旨} 金属の通常相とアモルファスの長距離秩序を比較する。

\Needspace{4\baselineskip}\paragraph{確認した条件・記述（独自要約）}\mbox{}\par

\noindent\textbf{判断対象：} 金属の通常相とアモルファスの長距離秩序を比較する。\par

\noindent\textbf{原典確認：} 全5候補・指数・否定・記号を保存原典22頁画像と照合。\par

\paragraph{問われる核心} 金属の通常相とアモルファスの長距離秩序を比較する。

\paragraph{必要知識} アモルファスは長距離の周期的結晶格子を持たないので5が不適切。

\paragraph{考え方・正答を選ぶ理由（査読待ち初稿）} アモルファスは長距離の周期的結晶格子を持たないので5が不適切。完全に全原子が無相関という意味ではない。Feの912℃の構造変態と770℃付近の磁気転移は別である。元素例は通常相を想定し全温度圧力に固定しない。

\Needspace{4\baselineskip}\paragraph{選択肢の検討（査読待ち初稿）}\mbox{}\par

\noindent\textbf{1：選ばない候補} Al/Cu/NiはFCC。通常の常温常圧相。\par

\noindent\textbf{2：選ばない候補} W/Mo/NbはBCC。通常の常温常圧相。\par

\noindent\textbf{3：選ばない候補} Mg/Zn/YはHCP。通常の六方相。\par

\noindent\textbf{4：選ばない候補} Feは912℃で格子型変態。純FeのαBCC→γFCCの相変態。\par

\noindent\textbf{5：選ぶ候補} アモルファスは規則的結晶格子。長距離の周期的並進秩序を持たない。短距離秩序はあり得る。\par

\paragraph{再利用する考え方} 軌道と電子状態、格子と基底、単位胞原子数と配位数、定義の分母を区別して候補を全て評価する。

\paragraph{共通理論} \hyperref[theory:atomic:geometry]{第\ref*{theory:atomic:geometry}節}

\par 参考：\cite{princeton_iron,princeton_tungsten,openstax_crystals}。

\paragraph{関連問題} \hyperref[q:r03-reg-applied-08]{R03-REG-APPLIED-08}

\par 難易度 2（校準初稿）、優先度未評価。

\section{電子配置・容量と原子データ}

\subsection{令和元年度再試験 問14}
\label{q:r01-ret-applied-14}

\noindent\textbf{原典：} sources/R01\_retest/17\_applied\_science.pdf、PDF物理5頁。

\textbf{公式正答：} 2（single）。

\paragraph{問題の要旨} 酸素基底状態の電子数とs軌道容量を確認する。

\Needspace{4\baselineskip}\paragraph{確認した条件・記述（独自要約）}\mbox{}\par

\noindent\textbf{判断対象：} 酸素基底状態の電子数とs軌道容量を確認する。\par

\noindent\textbf{原典確認：} 全5候補・指数・否定・記号を保存原典22頁画像と照合。\par

\paragraph{問われる核心} 酸素基底状態の電子数とs軌道容量を確認する。

\paragraph{必要知識} 酸素はZ=\allowbreak{}8で中性なら8電子。

\paragraph{考え方・正答を選ぶ理由（査読待ち初稿）} 酸素はZ=\allowbreak{}8で中性なら8電子。s副殻は1空間軌道に逆スピン2電子が最大で、基底配置1s2 2s2 2p4、2。候補は全て合計8でも、総数だけで判定できず容量とエネルギー順を確認する。

\Needspace{4\baselineskip}\paragraph{選択肢の検討（査読待ち初稿）}\mbox{}\par

\noindent\textbf{1：選ばない候補} 1s1 2s3 2p4。2sに3電子は不可能。\par

\noindent\textbf{2：選ぶ候補} 1s2 2s2 2p4。低い1s/2sを各2、2pを4で計8。\par

\noindent\textbf{3：選ばない候補} 1s1 2s4 2p3。2sに4電子は不可能。\par

\noindent\textbf{4：選ばない候補} 1s2 2s3 2p3。2sに3電子は不可能。\par

\noindent\textbf{5：選ばない候補} 1s2 2s4 2p2。2sに4電子は不可能。\par

\paragraph{再利用する考え方} 軌道と電子状態、格子と基底、単位胞原子数と配位数、定義の分母を区別して候補を全て評価する。

\paragraph{共通理論} \hyperref[theory:atomic:orbital]{第\ref*{theory:atomic:orbital}節}

\par 参考：\cite{atomic_quantum,atomic_nist2019}。

\paragraph{関連問題} \hyperref[q:r02-reg-applied-10]{R02-REG-APPLIED-10}

\par 難易度 2（校準初稿）、優先度未評価。

\subsection{令和2年度 問10}
\label{q:r02-reg-applied-10}

\noindent\textbf{原典：} sources/R02/17\_applied\_science.pdf、PDF物理4頁。

\textbf{公式正答：} ALL（all\_awarded）。

\par\textbf{独立解答と公式正答に相違あり：} 独立解答 3|5。公式ALLを維持。初回3|5は公式前の独立判断として保持。原典の銅・原子番号・電子数を別元素へ変更しない。公式措置理由の全文は未取得。

\paragraph{問題の要旨} 原典の複数の誤りと全員正解の公式採点を保持する。

\Needspace{4\baselineskip}\paragraph{確認した条件・記述（独自要約）}\mbox{}\par

\noindent\textbf{判断対象：} 原典の複数の誤りと全員正解の公式採点を保持する。\par

\noindent\textbf{原典確認：} 全5候補・指数・否定・記号を保存原典22頁画像と照合。\par

\paragraph{問われる核心} 原典の複数の誤りと全員正解の公式採点を保持する。

\paragraph{必要知識} 独立初回は3|5、公式はALL（全員正解）。

\paragraph{考え方・正答を選ぶ理由（査読待ち初稿）} 独立初回は3|5、公式はALL（全員正解）。原典には実際に銅とあり、Gaへ黙って修正しない。Cuは29、Znは30で4s電子も1と2。さらにKrの第一IEはArより小さいため5も不適切。単一候補へ絞らず公式採点と二つの科学的誤りを併記する。公式の全員正解措置の理由文は未取得なので、理由を公式発表として捏造しない。

\Needspace{4\baselineskip}\paragraph{選択肢の検討（査読待ち初稿）}\mbox{}\par

\noindent\textbf{1：公式採点は全員正解} Neは1s2 2s2 2p6。中性Neの10電子配置。\par

\noindent\textbf{2：公式採点は全員正解} Krは[Ar]4s2 3d10 4p6。中性Krの36電子配置。記載順と容量を区別。\par

\noindent\textbf{3：公式採点は全員正解} ZnよりZが1大きいCuは4sが1多い。原典のCuを保持。Cu29はZn30より1小さく4sは1対2で少ない。\par

\noindent\textbf{4：公式採点は全員正解} MgのIEはNaより大きい。NIST値Mg7.6462eV>Na5.1391eV。\par

\noindent\textbf{5：公式採点は全員正解} KrのIEはArより大きい。NIST値Kr13.9996eV<Ar15.7596eVで逆。\par

\paragraph{再利用する考え方} 軌道と電子状態、格子と基底、単位胞原子数と配位数、定義の分母を区別して候補を全て評価する。

\paragraph{共通理論} \hyperref[theory:atomic:orbital]{第\ref*{theory:atomic:orbital}節}

\par 参考：\cite{atomic_nist2019}。

\paragraph{関連問題} \hyperref[q:r01-ret-applied-14]{R01-RET-APPLIED-14}

\par 難易度 2（校準初稿）、優先度未評価。

\subsection{令和6年度 問8}
\label{q:r06-reg-applied-08}

\noindent\textbf{原典：} sources/R06/17\_applied\_science.pdf、PDF物理3頁。

\textbf{公式正答：} 4（single）。

\par\textbf{検討未確定：} 肢3の軌道という原典語は副殻群の総容量として条件を補う。肢5も孤立原子とバルク金属の常磁性機構が同じとはしない。

\paragraph{問題の要旨} 副殻容量・半導体の担い手とGeの四面体結合を区別する。

\Needspace{4\baselineskip}\paragraph{確認した条件・記述（独自要約）}\mbox{}\par

\noindent\textbf{原典確認：} S7\_051の11原典画像で13問の全五肢・指数・単位・否定を照合。\par

\noindent\textbf{対象：} 副殻容量・半導体の担い手とGeの四面体結合を区別する。\par

\paragraph{問われる核心} 副殻容量・半導体の担い手とGeの四面体結合を区別する。

\paragraph{必要知識} 4は四面体構造をsp\textsuperscript{2}へ置き換えたことが誤り。原子の4s\textsuperscript{2}4p\textsuperscript{2}配置と結晶の共有結合モデルを分ける。副殻の総容量と一個の空間軌道の容量、孤立原子と金属の磁気応答も分離する。

\paragraph{考え方・正答を選ぶ理由（査読待ち初稿）} 4は四面体構造をsp\textsuperscript{2}へ置き換えたことが誤り。原子の4s\textsuperscript{2}4p\textsuperscript{2}配置と結晶の共有結合モデルを分ける。副殻の総容量と一個の空間軌道の容量、孤立原子と金属の磁気応答も分離する。

\Needspace{4\baselineskip}\paragraph{選択肢の検討（査読待ち初稿）}\mbox{}\par

\noindent\textbf{1：代表的条件で適切} 電池で自由電子が正極側へ動く。電子の負電荷による平均ドリフトは通常の電流方向と反対。各電子の熱運動を含む瞬間的な全軌道が正極へ直線的に進むわけではない。\par

\noindent\textbf{2：代表的条件で適切} SiへBを添加するとp型。Bはアクセプターとして働き正孔が主要な担い手となる。結晶全体が単に正に帯電する意味ではない。\par

\noindent\textbf{3：条件・語義を留保} s・p・d・fの収容数2・6・10・14。同じ副殻に属する空間軌道の数1・3・5・7へ各2電子を入れた総容量。一個の空間軌道はどの種類も2電子なので原典の軌道という語を副殻群として説明する。\par

\noindent\textbf{4：不適切} Geの4s\textsuperscript{2}4p\textsuperscript{2}がsp\textsuperscript{2}でダイヤモンド構造。原典の電子配置は4s\textsuperscript{2}4p\textsuperscript{2}でNISTにも一致。通常のGeは四面体の4近接結合を持つダイヤモンド型で、教科書的混成モデルはsp\textsuperscript{3}。sp\textsuperscript{2}の三角平面と取り違え。\par

\noindent\textbf{5：条件・語義を留保} Li・Naの不対スピンと常磁性。孤立原子のs電子の軌道角運動量ゼロとスピンを区別する。ただしバルク金属のPauli応答はFermi統計を含み、孤立原子のCurie則をそのまま使わない。\par

\paragraph{再利用する考え方} 場・応答・担い手・材料・温度を区別し、数値の指数と単位、代表モデルの適用条件を確認する。

\paragraph{共通理論} \hyperref[theory:atomic:orbital]{第\ref*{theory:atomic:orbital}節}, \hyperref[theory:atomic:geometry]{第\ref*{theory:atomic:geometry}節}, \hyperref[theory:em:pauli]{第\ref*{theory:em:pauli}節}

\par 参考：\cite{atomic_quantum,em51_ge,em51_ge_structure,mit_semiconductor,em51_pauli}。

\par 難易度 3（校準初稿）、優先度未評価。

\section{分子間力・分散とLJのエネルギー}

\subsection{平成26年度 問11}
\label{q:h26-reg-applied-11}

\noindent\textbf{原典：} sources/H26/17\_applied\_science.pdf、PDF物理5頁。

\textbf{公式正答：} 2（single）。

\paragraph{問題の要旨} 分子極性と水素結合の供与・受容を判断する。

\Needspace{4\baselineskip}\paragraph{確認した条件・記述（独自要約）}\mbox{}\par

\noindent\textbf{判断対象：} 分子極性と水素結合の供与・受容を判断する。\par

\noindent\textbf{原典確認：} 全5候補・指数・否定・記号を保存原典22頁画像と照合。\par

\paragraph{問われる核心} 分子極性と水素結合の供与・受容を判断する。

\paragraph{必要知識} CH4は対称な正四面体で永久双極子がなく、極性が強いため水素結合を生じるという2が不適切。

\paragraph{考え方・正答を選ぶ理由（査読待ち初稿）} CH4は対称な正四面体で永久双極子がなく、極性が強いため水素結合を生じるという2が不適切。水・HFの典型的水素結合や異種分子間の供与体/受容体を区別する。水素を含めば全て同等の水素結合が生じるという推論はできない。

\Needspace{4\baselineskip}\paragraph{選択肢の検討（査読待ち初稿）}\mbox{}\par

\noindent\textbf{1：選ばない候補} 極性の強い水素化合物に多い。典型的なN/O/Fに結合したH等を入口とする。\par

\noindent\textbf{2：選ぶ候補} CH4は強極性で水素結合。正四面体CH4の結合双極子は相殺し、強極性という前提が誤り。\par

\noindent\textbf{3：選ばない候補} 水の沸点はH2Sより高い。水の分子間水素結合が高い沸点を説明する。\par

\noindent\textbf{4：選ばない候補} HFは高い電気陰性度に関係。HFの分極した結合と電子対による相互作用。\par

\noindent\textbf{5：選ばない候補} 異種分子間にも水素結合。供与体と受容体が違う分子でも成立する。\par

\paragraph{再利用する考え方} 軌道と電子状態、格子と基底、単位胞原子数と配位数、定義の分母を区別して候補を全て評価する。

\paragraph{共通理論} \hyperref[theory:atomic:interactions]{第\ref*{theory:atomic:interactions}節}

\par 参考：\cite{atomic_nonbond}。

\par 難易度 2（校準初稿）、優先度未評価。

\subsection{平成29年度 問11}
\label{q:h29-reg-applied-11}

\noindent\textbf{原典：} sources/H29/17\_applied\_science.pdf、PDF物理4頁。

\textbf{公式正答：} 5（single）。

\par\textbf{検討未確定：} 候補2は凝集エネルギーを配位数のみへ帰する排他的主張で疑義。候補5はポテンシャルと力の語も混在。公式5と科学的一意性を分けて査読。

\paragraph{問題の要旨} 結合多重度・分散力とLennard–Jonesの符号を確認する。

\Needspace{4\baselineskip}\paragraph{確認した条件・記述（独自要約）}\mbox{}\par

\noindent\textbf{判断対象：} 結合多重度・分散力とLennard–Jonesの符号を確認する。\par

\noindent\textbf{原典確認：} 全5候補・指数・否定・記号を保存原典22頁画像と照合。\par

\paragraph{問われる核心} 結合多重度・分散力とLennard–Jonesの符号を確認する。

\paragraph{必要知識} 通常U=\allowbreak{}4ε[(σ/r)\textasciicircum{}12−(σ/r)\textasciicircum{}6]で短距離反発が正のr\textasciicircum{}-12、引力が負のr\textasciicircum{}-6。

\paragraph{考え方・正答を選ぶ理由（査読待ち初稿）} 通常U=\allowbreak{}4ε[(σ/r)\textasciicircum{}12−(σ/r)\textasciicircum{}6]で短距離反発が正のr\textasciicircum{}-12、引力が負のr\textasciicircum{}-6。5は両者が逆なので独立初回も5。ただし力F=−dU/drでは指数13と7となり、原典の力という語も区別する。候補2のW凝集を配位数だけで説明する断定には別の疑義を残し、全他肢が科学的に無条件正しいとはしない。

\Needspace{4\baselineskip}\paragraph{選択肢の検討（査読待ち初稿）}\mbox{}\par

\noindent\textbf{1：選ばない候補} 炭素結合の多重度増で強く短くなる。代表的C-C/C=\allowbreak{}C/C≡Cの傾向。\par

\noindent\textbf{2：選ばない候補} Wの大凝集エネルギーは結合エネルギーでなく配位数だけ。配位数だけとの排他的説明に疑義。BCCの配位数8だけで凝集を保証できない。\par

\noindent\textbf{3：選ばない候補} CO2は電子対反発を避け直線。VSEPRでCまわり2電子領域は180度。\par

\noindent\textbf{4：選ばない候補} 分散は瞬間双極子による。瞬間双極子が他分子へ双極子を誘起。\par

\noindent\textbf{5：選ぶ候補} LJはr\textasciicircum{}-6反発とr\textasciicircum{}-12引力。通常の12–6ポテンシャルの反発と引力が逆。力の指数もポテンシャルとは異なる。\par

\paragraph{再利用する考え方} 軌道と電子状態、格子と基底、単位胞原子数と配位数、定義の分母を区別して候補を全て評価する。

\paragraph{共通理論} \hyperref[theory:atomic:interactions]{第\ref*{theory:atomic:interactions}節}

\par 参考：\cite{atomic_lj,atomic_nonbond}。

\par 難易度 2（校準初稿）、優先度未評価。

\section{NaCl/CsClと格子・基底}

\subsection{平成27年度 問9}
\label{q:h27-reg-applied-09}

\noindent\textbf{原典：} sources/H27/17\_applied\_science.pdf、PDF物理4頁。

\textbf{公式正答：} 3（single）。

\par\textbf{検討未確定：} 全隣接イオンが接するという記述はイオン半径の硬球近似。半径の定義・実結晶条件は原典で未指定。

\paragraph{問題の要旨} NaClの配位数と硬球近似の距離を区別する。

\Needspace{4\baselineskip}\paragraph{確認した条件・記述（独自要約）}\mbox{}\par

\noindent\textbf{判断対象：} NaClの配位数と硬球近似の距離を区別する。\par

\noindent\textbf{原典確認：} 全5候補・指数・否定・記号を保存原典22頁画像と照合。\par

\paragraph{問われる核心} NaClの配位数と硬球近似の距離を区別する。

\paragraph{必要知識} NaCl型は最近接異種イオンの配位数6で、4とする3が不適切。

\paragraph{考え方・正答を選ぶ理由（査読待ち初稿）} NaCl型は最近接異種イオンの配位数6で、4とする3が不適切。陽陰それぞれのFCC副格子がずれて重なる。硬球接触を置くとa/2=\allowbreak{}rNa+rCl、同種最近接a/√2=√2(rNa+rCl)。格子の幾何と、半径の和を実距離とするモデル仮定を分ける。

\Needspace{4\baselineskip}\paragraph{選択肢の検討（査読待ち初稿）}\mbox{}\par

\noindent\textbf{1：選ばない候補} 単位格子は立方体。慣用胞は立方。\par

\noindent\textbf{2：選ばない候補} 静電引力による結晶。陽陰イオンの静電相互作用。\par

\noindent\textbf{3：選ぶ候補} 最近接Clは4個。Naの周囲のClは上下左右前後の6個。\par

\noindent\textbf{4：選ばない候補} 隣接Na/Clは全て接触。イオン半径による硬球近似の記述で、電子雲の厳密境界ではない。\par

\noindent\textbf{5：選ばない候補} 最近接Na-Naは√2(rNa+rCl)。Na-Cl距離a/2=\allowbreak{}rNa+rClとNa-Na=\allowbreak{}a/√2から導く近似式。\par

\paragraph{再利用する考え方} 軌道と電子状態、格子と基底、単位胞原子数と配位数、定義の分母を区別して候補を全て評価する。

\paragraph{共通理論} \hyperref[theory:atomic:ionic]{第\ref*{theory:atomic:ionic}節}

\par 参考：\cite{openstax_crystals}。

\paragraph{関連問題} \hyperref[q:r03-reg-applied-14]{R03-REG-APPLIED-14}

\par 難易度 2（校準初稿）、優先度未評価。

\subsection{令和3年度 問14}
\label{q:r03-reg-applied-14}

\noindent\textbf{原典：} sources/R03/17\_applied\_science.pdf、PDF物理5頁。

\textbf{公式正答：} 3（single）。

\par\textbf{検討未確定：} 他候補は混成・典型結合強度・分子固体というモデル/一般傾向。全条件への無限定一般化は保持しない。

\paragraph{問題の要旨} NaClとCsClの格子・配位を比較する。

\Needspace{4\baselineskip}\paragraph{確認した条件・記述（独自要約）}\mbox{}\par

\noindent\textbf{判断対象：} NaClとCsClの格子・配位を比較する。\par

\noindent\textbf{原典確認：} 全5候補・指数・否定・記号を保存原典22頁画像と照合。\par

\paragraph{問われる核心} NaClとCsClの格子・配位を比較する。

\paragraph{必要知識} NaClは陽陰のFCC副格子を重ね配位6、CsClは単純立方の隅と異種の中心で配位8。

\paragraph{考え方・正答を選ぶ理由（査読待ち初稿）} NaClは陽陰のFCC副格子を重ね配位6、CsClは単純立方の隅と異種の中心で配位8。異なる幾何なので同構造/同Madelung定数という3は不適切。CsClをBCC Bravais格子と呼ばない。異種イオンの中心は隅と並進で同等ではない。Madelung定数の数値暗記を使わず構造の相違から判断する。

\Needspace{4\baselineskip}\paragraph{選択肢の検討（査読待ち初稿）}\mbox{}\par

\noindent\textbf{1：選ばない候補} O2/F2はσ分裂がπより大きい。指定2p軌道の重なりの向きの比較。\par

\noindent\textbf{2：選ばない候補} H2Oはsp3電子対で折れ線。結合対2と孤立対2のモデル。\par

\noindent\textbf{3：選ぶ候補} NaClとCsClは同構造・同Madelung定数。NaClはFCC基底で配位6、CsClは単純立方基底で配位8で異なる。\par

\noindent\textbf{4：選ばない候補} 水素結合は多くの場合共有結合より弱い。典型的分子間水素結合の比較で、全化学種の普遍的上限とはしない。\par

\noindent\textbf{5：選ばない候補} 分子性固体は一般に軟らかく低融点。相互作用の代表傾向で全分子固体の数値保証ではない。\par

\paragraph{再利用する考え方} 軌道と電子状態、格子と基底、単位胞原子数と配位数、定義の分母を区別して候補を全て評価する。

\paragraph{共通理論} \hyperref[theory:atomic:ionic]{第\ref*{theory:atomic:ionic}節}

\par 参考：\cite{openstax_crystals}。

\paragraph{関連問題} \hyperref[q:h27-reg-applied-09]{H27-REG-APPLIED-09}

\par 難易度 2（校準初稿）、優先度未評価。

\section{格子比熱・散乱と磁区応答}

\subsection{平成30年度 問9}
\label{q:h30-reg-applied-09}

\noindent\textbf{原典：} sources/H30/17\_applied\_science.pdf、PDF物理3頁。

\textbf{公式正答：} 5（single）。

\paragraph{問題の要旨} 格子振動の比熱・散乱と磁区の応答を区別する。

\Needspace{4\baselineskip}\paragraph{確認した条件・記述（独自要約）}\mbox{}\par

\noindent\textbf{判断対象：} 格子振動の比熱・散乱と磁区の応答を区別する。\par

\noindent\textbf{原典確認：} 全5候補・指数・否定・記号を保存原典22頁画像と照合。\par

\paragraph{問われる核心} 格子振動の比熱・散乱と磁区の応答を区別する。

\paragraph{必要知識} 1から4は格子の振動自由度又は熱変位を直接扱う。

\paragraph{考え方・正答を選ぶ理由（査読待ち初稿）} 1から4は格子の振動自由度又は熱変位を直接扱う。Barkhausen効果は磁区壁がピン止めを越えて不連続に動く磁化応答で5。磁性と格子が一切相互作用しないという断言へ広げない。高温3Rと低温で減る格子比熱は振動子の量子化の同じ枠組みで説明できる。

\Needspace{4\baselineskip}\paragraph{選択肢の検討（査読待ち初稿）}\mbox{}\par

\noindent\textbf{1：選ばない候補} Dulong–Petit則。高温の格子比熱3Rに対応。\par

\noindent\textbf{2：選ばない候補} Debye–Waller因子。原子の変位による散乱強度の減衰。\par

\noindent\textbf{3：選ばない候補} Debye温度。格子振動スペクトルの特徴温度。\par

\noindent\textbf{4：選ばない候補} Einstein比熱。量子調和振動子による格子比熱モデル。\par

\noindent\textbf{5：選ぶ候補} Barkhausen効果。磁区壁の不連続運動による磁化変化で、比較中は格子振動との関係が薄い。\par

\paragraph{再利用する考え方} 軌道と電子状態、格子と基底、単位胞原子数と配位数、定義の分母を区別して候補を全て評価する。

\paragraph{共通理論} \hyperref[theory:atomic:phonon]{第\ref*{theory:atomic:phonon}節}

\par 参考：\cite{atomic_phonon,atomic_barkhausen}。

\par 難易度 2（校準初稿）、優先度未評価。

\section{空間軌道・スピンと結合対称性}

\subsection{令和元年度 問14}
\label{q:r01-reg-applied-14}

\noindent\textbf{原典：} sources/R01/17\_applied\_science.pdf、PDF物理4頁。

\textbf{公式正答：} 1（single）。

\par\textbf{独立解答と公式正答に相違あり：} 独立解答 4。初回4/公式1の不一致。φ相互作用と純φ結合の分類、水素の空間軌道と電子状態、d混成近似の解釈を要査読。公式に合わせて不存在の説明を作らない。

\paragraph{問題の要旨} 空間軌道・スピンと結合対称性を区別する。

\Needspace{4\baselineskip}\paragraph{確認した条件・記述（独自要約）}\mbox{}\par

\noindent\textbf{判断対象：} 空間軌道・スピンと結合対称性を区別する。\par

\noindent\textbf{原典確認：} 全5候補・指数・否定・記号を保存原典22頁画像と照合。\par

\paragraph{問われる核心} 空間軌道・スピンと結合対称性を区別する。

\paragraph{必要知識} 独立初回は4、公式は1で不一致を保持。

\paragraph{考え方・正答を選ぶ理由（査読待ち初稿）} 独立初回は4、公式は1で不一致を保持。n,l,mは空間軌道、スピンは電子状態の属性で水素にも存在する。候補4を多電子にのみスピンがあると読むと問題がある。一方、σ/π/δ/φは軸まわりの軌道対称性に対応し、2014年のthorocene研究にはφ型軌道相互作用の実験的根拠がある。これを単純な二原子純φ結合の確立と同一視もしない。公式1の採点を表示しつつ、1の一意の不適切性は未解決としてAstraへ送る。

\Needspace{4\baselineskip}\paragraph{選択肢の検討（査読待ち初稿）}\mbox{}\par

\noindent\textbf{1：選ぶ候補} σ/π/δ/φ等の結合。公式はこの候補1。ただし2014年研究にはφ軌道相互作用があり、φの概念の不存在では説明できない。\par

\noindent\textbf{2：選ばない候補} H2OはAB2E2。O周囲は結合対2と孤立対2。\par

\noindent\textbf{3：選ばない候補} dを含む混成もある。d関与の形式的混成モデルはある。超原子価をd混成だけで実在説明するのは別問題。\par

\noindent\textbf{4：選ばない候補} 水素はn,l,mだが多電子では加えてスピン。空間軌道は3量子数、電子状態は水素にもスピンがある。多電子にのみ追加する読みには疑義。\par

\noindent\textbf{5：選ばない候補} 3p充填後は4s。中性原子の通常の構築順でAr後Kに4sが入る。全励起状態・イオンの順序とはしない。\par

\paragraph{再利用する考え方} 軌道と電子状態、格子と基底、単位胞原子数と配位数、定義の分母を区別して候補を全て評価する。

\paragraph{共通理論} \hyperref[theory:atomic:orbital]{第\ref*{theory:atomic:orbital}節}

\par 参考：\cite{atomic_quantum,atomic_phi2014}。

\par 難易度 2（校準初稿）、優先度未評価。

\section{炭素同素体の構造と物性}

\subsection{平成28年度 問14}
\label{q:h28-reg-applied-14}

\noindent\textbf{原典：} sources/H28/17\_applied\_science.pdf、PDF物理5頁。

\textbf{公式正答：} 2（single）。

\par\textbf{検討未確定：} ダイヤモンドの硬度をMohs尺度・通常純結晶へ限定。ドーピング等の導電性例外を保持。

\paragraph{問題の要旨} 炭素同素体の結合構造と導電性を比較する。

\Needspace{4\baselineskip}\paragraph{確認した条件・記述（独自要約）}\mbox{}\par

\noindent\textbf{判断対象：} 炭素同素体の結合構造と導電性を比較する。\par

\noindent\textbf{原典確認：} 全5候補・指数・否定・記号を保存原典22頁画像と照合。\par

\paragraph{問われる核心} 炭素同素体の結合構造と導電性を比較する。

\paragraph{必要知識} 通常の純粋ダイヤモンドはsp3結合網で硬いが良導体ではない。

\paragraph{考え方・正答を選ぶ理由（査読待ち初稿）} 通常の純粋ダイヤモンドはsp3結合網で硬いが良導体ではない。共有結合なら導電性が高いとする候補2が不適切。黒鉛の層内/層間と、CNTの巻き方依存の電子構造を区別する。ホウ素ドープによる導電・超伝導ダイヤモンドの例外は通常純結晶の説明と別に扱う。

\Needspace{4\baselineskip}\paragraph{選択肢の検討（査読待ち初稿）}\mbox{}\par

\noindent\textbf{1：選ばない候補} ダイヤモンドは四面体結合で全鉱物中最高硬度。通常の鉱物のMohs尺度の最高位10。全人工材料・全荷重の絶対比較とはしない。\par

\noindent\textbf{2：選ぶ候補} ダイヤモンドは共有結合結晶で電気をよく通す。通常の純粋ダイヤモンドは大きなバンドギャップを持ち良導体ではない。ドープ材の例外と分ける。\par

\noindent\textbf{3：選ばない候補} 黒鉛は六角網面で層間が弱い。層内の強い共有結合と層間相互作用を分ける。\par

\noindent\textbf{4：選ばない候補} CNTは単層又は多層同軸のナノ管。六角網面の管状構造。原典の同軸形状を保持。\par

\noindent\textbf{5：選ばない候補} CNTの直径・配列で導電性が変わる。巻き方のカイラリティと径が電子構造へ関係する。\par

\paragraph{再利用する考え方} 軌道と電子状態、格子と基底、単位胞原子数と配位数、定義の分母を区別して候補を全て評価する。

\paragraph{共通理論} \hyperref[theory:atomic:carbon]{第\ref*{theory:atomic:carbon}節}

\par 参考：\cite{atomic_cnt,atomic_diamond}。

\paragraph{関連問題} \hyperref[q:r05-reg-applied-10]{R05-REG-APPLIED-10}

\par 難易度 2（校準初稿）、優先度未評価。

\subsection{令和5年度 問10}
\label{q:r05-reg-applied-10}

\noindent\textbf{原典：} sources/R05/17\_applied\_science.pdf、PDF物理4頁。

\textbf{公式正答：} 5（single）。

\par\textbf{検討未確定：} ダイヤモンドの硬度をMohs尺度・通常純結晶へ限定。ドーピング等の導電性例外を保持。

\paragraph{問題の要旨} 炭素同素体の結合構造と導電性を比較する。

\Needspace{4\baselineskip}\paragraph{確認した条件・記述（独自要約）}\mbox{}\par

\noindent\textbf{判断対象：} 炭素同素体の結合構造と導電性を比較する。\par

\noindent\textbf{原典確認：} 全5候補・指数・否定・記号を保存原典22頁画像と照合。\par

\paragraph{問われる核心} 炭素同素体の結合構造と導電性を比較する。

\paragraph{必要知識} 通常の純粋ダイヤモンドはsp3結合網で硬いが良導体ではない。

\paragraph{考え方・正答を選ぶ理由（査読待ち初稿）} 通常の純粋ダイヤモンドはsp3結合網で硬いが良導体ではない。共有結合なら導電性が高いとする候補5が不適切。黒鉛の層内/層間と、CNTの巻き方依存の電子構造を区別する。ホウ素ドープによる導電・超伝導ダイヤモンドの例外は通常純結晶の説明と別に扱う。

\Needspace{4\baselineskip}\paragraph{選択肢の検討（査読待ち初稿）}\mbox{}\par

\noindent\textbf{1：選ばない候補} CNTは単層又は多層同軸のナノ管。六角網面の管状構造。原典の同軸形状を保持。\par

\noindent\textbf{2：選ばない候補} CNTの直径・配列で導電性が変わる。巻き方のカイラリティと径が電子構造へ関係する。\par

\noindent\textbf{3：選ばない候補} 黒鉛は六角網面で層間が弱い。層内の強い共有結合と層間相互作用を分ける。\par

\noindent\textbf{4：選ばない候補} ダイヤモンドは四面体結合で全鉱物中最高硬度。通常の鉱物のMohs尺度の最高位10。全人工材料・全荷重の絶対比較とはしない。\par

\noindent\textbf{5：選ぶ候補} ダイヤモンドは共有結合結晶で電気をよく通す。通常の純粋ダイヤモンドは大きなバンドギャップを持ち良導体ではない。ドープ材の例外と分ける。\par

\paragraph{再利用する考え方} 軌道と電子状態、格子と基底、単位胞原子数と配位数、定義の分母を区別して候補を全て評価する。

\paragraph{共通理論} \hyperref[theory:atomic:carbon]{第\ref*{theory:atomic:carbon}節}

\par 参考：\cite{atomic_cnt,atomic_diamond}。

\paragraph{関連問題} \hyperref[q:h28-reg-applied-14]{H28-REG-APPLIED-14}

\par 難易度 2（校準初稿）、優先度未評価。

\section{同位体の基礎}

\subsection{平成28年度 問10}
\label{q:h28-reg-applied-10}

\noindent\textbf{原典：} sources/H28/17\_applied\_science.pdf、PDF物理4頁。

\textbf{公式正答：} 2（single）。

\paragraph{問題の要旨} 原子番号と質量数から同位体を定義する。

\Needspace{4\baselineskip}\paragraph{確認した条件・記述（独自要約）}\mbox{}\par

\noindent\textbf{判断対象：} 原子番号と質量数から同位体を定義する。\par

\noindent\textbf{原典確認：} 全5候補・指数・否定・記号を保存原典22頁画像と照合。\par

\paragraph{問われる核心} 原子番号と質量数から同位体を定義する。

\paragraph{必要知識} 同位体はZが等しくNが違う原子の組で、A=\allowbreak{}Z+Nは異なる。

\paragraph{考え方・正答を選ぶ理由（査読待ち初稿）} 同位体はZが等しくNが違う原子の組で、A=\allowbreak{}Z+Nは異なる。2。1は同素体、3は同重体、4は周期表の族、5は有機同族列。単体の結晶構造が違うことと核の中性子数が違うことを混同しない。

\Needspace{4\baselineskip}\paragraph{選択肢の検討（査読待ち初稿）}\mbox{}\par

\noindent\textbf{1：選ばない候補} 硫黄の異なる単体。同素体の説明。\par

\noindent\textbf{2：選ぶ候補} 原子番号同一・質量数不同。同じZ、違う中性子数N、A=\allowbreak{}Z+Nが異なる。\par

\noindent\textbf{3：選ばない候補} 陽子と中性子の和が同じ。同じAの同重体でZ同一とは限らない。\par

\noindent\textbf{4：選ばない候補} 周期表の同族。電子配置による化学的族の分類。\par

\noindent\textbf{5：選ばない候補} アルカン等の同族列。有機化合物の同族列で核種ではない。\par

\paragraph{再利用する考え方} 軌道と電子状態、格子と基底、単位胞原子数と配位数、定義の分母を区別して候補を全て評価する。

\paragraph{共通理論} \hyperref[theory:atomic:isotopes]{第\ref*{theory:atomic:isotopes}節}

\par 参考：\cite{modern_mit_nuclear}。

\paragraph{関連問題} \hyperref[q:h29-reg-applied-18]{H29-REG-APPLIED-18}

\par 難易度 2（校準初稿）、優先度未評価。

